% Generated draft: requires mathematical and source review.
\documentclass[11pt]{article}
\usepackage[margin=25mm]{geometry}
\usepackage{fontspec}
\setmainfont{Latin Modern Roman}
\usepackage{amsmath,amssymb,mathtools,unicode-math}
\setmathfont{Latin Modern Math}
\usepackage{enumitem,needspace,xurl,hyperref}
\setlength{\parindent}{0pt}
\setlength{\parskip}{5pt}
\setlength{\emergencystretch}{3em}
\newcommand{\N}{\mathbb{N}}
\newcommand{\Z}{\mathbb{Z}}
\newcommand{\Q}{\mathbb{Q}}
\newcommand{\R}{\mathbb{R}}
\newcommand{\C}{\mathbb{C}}
\newcommand{\F}{\mathbb{F}}
\newcommand{\A}{\mathbb{A}}
\newcommand{\PP}{\mathbb{P}}
\newcommand{\E}{\mathbb{E}}
\newcommand{\eps}{\varepsilon}
\newcommand{\dd}{\,\mathrm{d}}
\begin{document}
\section*{UnsolvedMath ID: 20002018}
\textbf{Source problem code:} catalogue entry
\subsection{Definitions and mathematical statement}
Problem 1b asks how to extract information by varying a metric across conformal classes. Its antecedent, companion record 20002017, concerns functionals well understood under conformal changes and contrasts local curvature integrals with nonlocal spectral invariants such as functional determinants. It does not specify a unique functional, dimension or desired invariant, and does not assume every such functional is conformally invariant.

A precise first-order model is the following. On a closed manifold let \(\mathcal F\) be differentiable, diffeomorphism invariant and conformally invariant, and assume it has an \(L^2\) gradient \(E_g\):
\[
D\mathcal F_g(h)=\int_M\langle E_g,h\rangle_g\,\dd V_g.
\]
Conformal and diffeomorphism invariance imply \(\operatorname{tr}_gE_g=0\) and \(\operatorname{div}_gE_g=0\). Put
\[
\mathcal T_g=\{h\in\Gamma(\operatorname{Sym}^2T^*M):\operatorname{tr}_gh=0,\ \operatorname{div}_gh=0\}.
\]
The aid reports that the maximal unit transverse first-order signal is
\[
\sup_{h\in\mathcal T_g,\,\|h\|_{L^2}=1}|D\mathcal F_g(h)|=\|E_g\|_{L^2},
\]
with value zero when the gradient vanishes. This uses the invariant functional and existence of a genuine \(L^2\) gradient; knowledge of a conformal anomaly alone does not give these hypotheses.

In real dimension four, the aid specializes to the Weyl-curvature functional \(\mathcal W\). Use the LeBrun normalization stated in the record, defining the Bach tensor's coefficient/sign by
\[
D\mathcal W_g(h)=-\int_M\langle B_g,h\rangle_g\,\dd V_g.
\]
Here \(B_g\) is trace free and divergence free. Its norm \(S_{\mathcal W}(g)=\|B_g\|_{L^2}\) is the reported maximal transverse signal. Bach-flatness \(B_g=0\) is equivalent to stationarity under all trace-free variations, and also under all transverse-traceless variations. If \(B_g\ne0\), the direction \(h=B_g\) gives
\[
D\mathcal W_g(B_g)=-\|B_g\|_{L^2}^2<0.
\]
The normalization of the Weyl norm and Bach tensor must match this variational formula; an alternative convention can change numerical factors.

The aid also defines
\[
X_{\mathcal W}(g)=\operatorname{Vol}_g(M)\int_M|B_g|_g^2\,\dd V_g.
\]
It is invariant under constant rescaling \(g\mapsto c^2g\), \(c>0\). This is not a claim of invariance under arbitrary point-dependent conformal changes. Vanishing of the first-order signal means stationarity at that metric, not absence of higher-order information or a complete answer for nonlocal functionals. Thus the certificate is a source-focused partial method, not an exhaustive solution of the broad question.
\subsection{Short English statement}
The source asks how to detect changes transverse to conformal classes. The supplied aid offers a first-order variational certificate: for the four-dimensional Weyl functional the Bach tensor measures the maximal transverse signal. It does not compute arbitrary nonlocal spectral variations.
\subsection{Sources}
\begin{itemize}
\item Thomas Branson, \emph{Anti-conformal perturbations}, Problem 1b, AIM workshop \emph{Conformal structure in geometry, analysis, and physics}, 12--16 August 2003: \url{https://aimath.org/WWN/confstruct/confstruct.pdf}. Attribution and functional context are supplied by adjacent Problem 1a.
\item Local registry \texttt{data/unsolvedmath/problems.json}, permanent ID 20002018, code \texttt{catalogue entry}; canonical \texttt{aim-geometry-notes.json}, \texttt{notes[355]}; updated 14 August 2026. Companion ID 20002017, code \texttt{AIM-GEOMETRY-0355}, \texttt{notes[354]}, gives Problem 1a and its local-versus-spectral distinction.
\item The aid names LeBrun's convention without a full bibliographic reference; the original literature field is empty. Attempt 1: partial result, partially solved status, rigor label proved, novelty confidence low. It is machine generated, not peer reviewed, and requires independent verification before citation.
\end{itemize}
\end{document}
